\chapter{1976 Nobel Prize in Economics}
\textbf{Laureates: }Milton Friedman (USA)

\section*{Reason for Award}
For his achievements in the fields of consumption analysis, monetary history and theory, and for
his demonstration of the complexity of stabilization policy

\section{What They Did}
Milton Friedman championed free-market capitalism and developed the permanent income hypothesis,
proposing that consumption depends on expected lifetime income rather than current income. This
hypothesis explained why temporary tax changes have limited effects on consumption, while
permanent changes have larger effects.

He provided a historical analysis of money's role in economic fluctuations, challenging Keynesian
orthodoxy by emphasizing the quantity theory of money. Friedman demonstrated that inflation is
always and everywhere a monetary phenomenon, requiring stable monetary growth rules rather than
discretionary policy. His analysis of the Great Depression showed that it resulted from monetary
contraction rather than market failure.

Friedman demonstrated that the Phillips curve trade-off between inflation and unemployment exists
only temporarily, as expectations adjust. He developed the theory of the natural rate of
unemployment, arguing that attempts to push unemployment below this rate merely accelerate inflation. Friedman also contributed to the theory of controlled experimentation, natural
experiments, and the measurement of money demand.

His advocacy for flexible exchange rates influenced currency regime debates worldwide. He developed
the concept of adaptive expectations, showing how people form expectations based on past
experiences. Friedman's work on the consumption function revolutionized macroeconomic analysis by
emphasizing forward-looking behavior.

Friedman also made important contributions to the theory of fiscal policy, analyzing the effects
of government spending and taxation on aggregate demand. His work on the money demand function
provided a theoretical foundation for monetarist policy recommendations. Friedman's "A Monetary
History of the United States" (1963), co-authored with Anna Schwartz, provided a comprehensive
analysis of monetary events in American history.

Friedman's work on the monetary history of the United States demonstrated the importance of
monetary factors in explaining economic fluctuations, challenging the prevailing Keynesian view
that fiscal policy was the primary tool for stabilization.

\section{Impact on Economic Thinking}
Friedman reinvigorated classical monetary theory, arguing that monetary policy is the primary
determinant of nominal economic fluctuations. He rejected active fiscal stabilization in favor of
rules-based policy, advocating a single rule for monetary growth that would provide stability and
predictability.

His monetarist critique inspired deregulation movements and skepticism about government
intervention, emphasizing individual freedom and market solutions. The permanent income hypothesis
reshaped consumption-savings analysis, linking macroeconomic stability to microeconomic decisions
about lifetime resource allocation.

Friedman's work on the natural rate of unemployment clarified the limits of monetary policy,
showing that attempts to push unemployment below its natural rate merely accelerate inflation. His
analysis of the Great Depression as a monetary phenomenon changed how economists understand
financial crises and the importance of central bank responsibility.

Friedman's policy recommendations influenced the shift toward rules-based monetary policy and
inflation targeting that became standard in many countries. His critique of discretionary
stabilization policy highlighted the time inconsistency problem and the importance of credible
commitments. The expectation-augmented Phillips curve became central to macroeconomic theory,
incorporating Friedman's insights about the role of expectations.

His work on the permanent income hypothesis influenced consumer theory and the analysis of
saving behavior. Friedman's advocacy for school vouchers and privatization influenced debates
about education policy and government provision of public services. His ideas about the role of
expectations in economic decision-making have had lasting influence on macroeconomic theory.

Friedman's influence extended beyond academic economics to public policy debates, where his
arguments for free markets and limited government intervention shaped political discourse in
many countries.

\section{Modern Perspectives Today}
Monetarist ideas continue influencing central banking through adopted inflation targeting
frameworks that reflect Friedman's critiques of discretionary policy. His arguments against
fine-tuning resonate in debates about quantitative easing and unconventional monetary tools.

The permanent income hypothesis remains the foundation of modern consumption theory, informing
models of saving, borrowing, and wealth accumulation. Friedman influenced supply-side economics
and Reagan-Thatcher era policies through his emphasis on incentives and market mechanisms. His
advocacy for flexible exchange rates influenced currency regime choices worldwide.

School choice vouchers and privatization debates draw on his institutional analysis. The ongoing
tension between rules and discretion in monetary policy reflects Friedman's foundational insights.
His legacy persists in how economists think about the limits of government intervention, the
importance of expectations, and the fundamental role of monetary stability in economic performance.

The 2008 financial crisis prompted renewed interest in Friedman's warnings about the dangers of
monetary contraction and the importance of lender of last resort functions. Modern central banking
practice incorporates many of Friedman's insights about the importance of credibility and
rules-based approaches to monetary policy. His work on the monetary history of the United States
remains a classic reference for understanding the role of monetary policy in economic fluctuations.

Contemporary debates about monetary policy rules versus discretion, the role of central banks,
and the importance of credible commitments continue to engage with Friedman's insights. His
work on the expectations-augmented Phillips curve remains central to modern macroeconomic theory
and policy analysis.
